% In this file you should put the actual content of the blueprint.
% It will be used both by the web and the print version.
% It should *not* include the \begin{document}
%
% If you want to split the blueprint content into several files then
% the current file can be a simple sequence of \input. Otherwise It
% can start with a \section or \chapter for instance.

\chapter{Setup}

\section{The cyclotomic field $K$}

Let $p$ be a prime number and $f$ a positive integer. Set $q = p^f$.
Let $K = \QQ(\zeta_{q-1})$ be the $(q-1)$-th cyclotomic field.
We fix a maximal ideal $P$ of $\cO_K$ lying above $p$.

\section{The Teichmüller character}

\begin{definition}
  \label{def:teichmuller}
  \lean{teichmuller}
  \leanok
  The \emph{Teichmüller character} is the unique group isomorphism
  \[
    \om : (\cO_K/P)^\times \xrightarrow{\sim} \mu_{q-1}(\cO_K)
  \]
  satisfying $\om(x) \equiv x \pmod{P}$ for all $x \in (\cO_K/P)^\times$,
  extended to all of $\cO_K/P$ by $\om(0) = 0$.
\end{definition}

\begin{lemma}
  \label{lem:orderOf_teichmuller}
  \uses{def:teichmuller}
  \lean{orderOf_teichmuller}
  \leanok
  The order of $\om$ as a multiplicative character is $q - 1$.
\end{lemma}

\begin{proof}
  \leanok
  Since $\om$ is an isomorphism between groups of order $q - 1$, it has
  order $q - 1$ as a group homomorphism.
\end{proof}

\begin{lemma}
  \label{lem:map_teichmuller_zpow_eq}
  \lean{map_teichmuller_zpow_eq}
  \leanok
  \uses{def:teichmuller}
  Let $R$ be a ring containing $\cO_K$ and let $\sigma$ be a ring endomorphism of $R$ such that
  $\sigma(\zeta_{q-1}) = \zeta_{q-1}^m$ for some $m \in \ZZ$. Then
  $\sigma(\om(x)) = \om(x)^m$ for all $x \in \cO_K/P$.
\end{lemma}

\begin{proof}
  \leanok
  Since $\om(x) \in \mu_{q-1}(\cO_K)$, we have $\om(x) = \zeta_{q-1}^k$
  for some $k$. Then $\sigma(\om(x)) = \sigma(\zeta_{q-1})^k = \zeta_{q-1}^{mk}
  = \om(x)^m$.
\end{proof}

\begin{lemma}
  \label{lem:teichmuller_sum}
  \uses{def:teichmuller, lem:orderOf_teichmuller}
  \leanok
  For all $a \in \ZZ$ with $(q-1) \nmid a$,
  \[
    \sum_{x \in \cO_K/P} \om(x)^{-a} = 0.
  \]
\end{lemma}

\begin{proof}
  \leanok
  Since $(q-1) \nmid a$, the character $\om^{-a}$ is nontrivial by
  \ref{lem:orderOf_teichmuller}. By \texttt{MulChar.sum\_eq\_zero\_of\_ne\_one},
  $\sum_{x \in \cO_K/P} \om(x)^{-a} = 0$.
\end{proof}

\section{The additive character}

Let $L$ be an extension of $K$ containing a primitive $p$-th root of unity
$\zeta_p$. We fix such a $\zeta_p \in \cO_L$.

\begin{definition}
  \label{def:addchar}
  \lean{addCharTrace}
  \leanok
  The \emph{additive character} $\psi : \cO_K/P \to \mu_p(\cO_L)$ is defined by
  \[
    \psi(x) = \zeta_p^{\tr_{(\cO_K/P)/(\ZZ/p\ZZ)}(x)}
  \]
  where the exponent is lifted canonically to $\ZZ$.
\end{definition}

\begin{lemma}
  \label{lem:addCharTrace_ne_one}
  \uses{def:addchar}
  \lean{addCharTrace_ne_one}
  \leanok
  The character $\psi$ is nontrivial, i.e., $\psi \neq 1$.
\end{lemma}

\begin{proof}
  \leanok
  Since the extension $\cO_K/P$ over $\ZZ/p\ZZ$ is separable, the trace
  $\tr : \cO_K/P \to \ZZ/p\ZZ$ is nonzero, so there exists $x$ with
  $\tr(x) \neq 0$, hence $\psi(x) = \zeta_p^{\tr(x)} \neq 1$.
\end{proof}

\begin{lemma}
  \label{lem:addCharTrace_isPrimitive}
  \uses{lem:addCharTrace_ne_one}
  \lean{addCharTrace_isPrimitive}
  \leanok
  The character $\psi$ is primitive.
\end{lemma}

\begin{proof}
  \leanok
  Since $\cO_K/P$ is a field, any nontrivial additive character is primitive.
  The character $\psi$ is nontrivial by \ref{lem:addCharTrace_ne_one}.
\end{proof}

\begin{lemma}
  \label{lem:addCharTrace_frob_apply}
  \uses{def:addchar}
  \lean{addCharTrace_frob_apply}
  \leanok
  For all $x \in \cO_K/P$, $\psi(x^p) = \psi(x)$.
\end{lemma}

\begin{proof}
  \leanok
  The Frobenius map $\varphi : x \mapsto x^p$ is a $(\ZZ/p\ZZ)$-algebra
  automorphism of $\cO_K/P$. By \texttt{Algebra.trace\_eq\_of\_algEquiv},
  $\tr(\varphi(x)) = \tr(x)$, i.e., $\tr(x^p) = \tr(x)$. Therefore
  $\psi(x^p) = \zeta_p^{\tr(x^p)} = \zeta_p^{\tr(x)} = \psi(x)$.
\end{proof}

\begin{lemma}
  \label{lem:addCharTrace_mk_sq_eq}
  \uses{def:addchar}
  \lean{addCharTrace_mk_sq_eq}
  \leanok
  Let $\fP$ be a prime ideal of $\cO_L$ with $\zeta_p - 1 \in \fP$.
  Then
  \[
    \psi(x) \equiv 1 + \tr(x) \cdot (\zeta_p - 1) \pmod{\fP^2}
  \]
  for all $x \in \cO_K/P$.
\end{lemma}

\begin{proof}
  \leanok
  There exists $a \in \mathbb{N}$ such that $\tr(x) = a$ and $\psi(x) = \zeta_p^a$.
  By the binomial theorem,
  \[
    \psi(x) = \zeta_p^a = (1 + (\zeta_p - 1))^a = \sum_{n=0}^{a} \binom{a}{n} (\zeta_p - 1)^n.
  \]
  Since $\zeta_p - 1 \in \fP$, we have $(\zeta_p - 1)^n \in \fP^n \subseteq \fP^2$ for all
  $n \geq 2$. Therefore
  \[
    \psi(x) \equiv 1 + a(\zeta_p - 1) = 1 + \tr(x) \cdot (\zeta_p - 1) \pmod{\fP^2}.
  \]
\end{proof}

\begin{lemma}
  \label{lem:addCharTrace_mk_eq_one}
  \uses{lem:addCharTrace_mk_sq_eq}
  \lean{addCharTrace_mk_eq_one}
  \leanok
  Let $\fP$ be a prime ideal of $\cO_L$ with $\zeta_p - 1 \in \fP$.
  Then $\psi(x) \equiv 1 \pmod{\fP}$ for all $x \in \cO_K/P$.
\end{lemma}

\begin{proof}
  \leanok
  By \ref{lem:addCharTrace_mk_sq_eq}, $\psi(x) \equiv 1 + \tr(x) \cdot (\zeta_p - 1) \pmod{\fP^2}$.
  Since $\zeta_p - 1 \in \fP$, the term $\tr(x) \cdot (\zeta_p - 1) \in \fP$,
  hence $\psi(x) \equiv 1 \pmod{\fP}$.
\end{proof}

\chapter{The Gauss sum}
\label{chap:gauss_sum}

Throughout this chapter and the next two, we assume that $p$ is an odd prime.

Let $\fP$ be a prime ideal of $\cO_L$ with $\zeta_p - 1 \in \fP$. Let $F$ be
the subfield of $L$ with $F = \QQ(\zeta_p)$. The Galois group $\Gal(L/F)$
is isomorphic to $(\ZZ/(q-1)\ZZ)^\times$ via the map $\sigma \mapsto a$
where $\sigma(\zeta_{q-1}) = \zeta_{q-1}^a$. We denote by $\fn(\sigma) \in
\{1, \ldots, q-2\}$ the natural number corresponding to $\sigma$ under this
isomorphism.

\section{Gauss sums}

\begin{definition}
  \label{def:gauss_sum_s}
  \uses{def:teichmuller, def:addchar}
  \lean{GaussSum}
  \leanok
  For $a \in \ZZ$, the \emph{Gauss sum} is
  \[
    g_a = g(\om^{-a}, \psi) = \sum_{x \in \cO_K/P} \om(x)^{-a} \cdot \psi(x).
  \]
\end{definition}

\begin{lemma}
  \label{lem:GaussSum_periodic}
  \uses{def:gauss_sum_s, lem:orderOf_teichmuller}
  \lean{GaussSum_periodic}
  \leanok
  If $(q-1) \mid k$, then for all $a \in \ZZ$, $g_{a+k} = g_a$.
\end{lemma}

\begin{proof}
  \leanok
  Since $\omega^{q-1} = 1$ by \ref{lem:orderOf_teichmuller}, $(q-1) \mid k$ implies
  $\omega^{-k} = 1$, so $\omega^{-(a+k)} = \omega^{-a}$, hence $g_{a+k} = g_a$.
\end{proof}

\begin{lemma}
  \label{lem:GaussSum_mem}
  \uses{def:gauss_sum_s, lem:addCharTrace_mk_eq_one, lem:orderOf_teichmuller, lem:teichmuller_sum}
  \lean{GaussSum_mem}
  \leanok
  If $(q-1) \nmid a$, then $g_a \in \fP$.
\end{lemma}

\begin{proof}
  \leanok
  By \ref{lem:addCharTrace_mk_eq_one}, $\psi(x) \equiv 1 \pmod{\fP}$ for all $x \in \cO_K/P$.
  Therefore
  \[
    g_a = \sum_{x \in \cO_K/P} \omega(x)^{-a} \psi(x)
    \equiv \sum_{x \in \cO_K/P} \omega(x)^{-a} \pmod{\fP}.
  \]
  Since $(q-1) \nmid a$, we have $\omega^{-a} \neq 1$ by \ref{lem:orderOf_teichmuller}.
  Hence $\sum_{x \in \cO_K/P} \omega(x)^{-a} = 0$ by \ref{lem:teichmuller_sum}.
  Therefore $g_a \equiv 0 \pmod{\fP}$, i.e.\ $g_a \in \fP$.
\end{proof}

\begin{lemma}
  \label{lem:GaussSum_mul_GaussSum_neg}
  \uses{def:gauss_sum_s, lem:addCharTrace_isPrimitive, def:teichmuller}
  \lean{GaussSum_mul_GaussSum_neg}
  \leanok
  If $(q-1) \nmid a$, then
  \[
    g_a \cdot g_{-a} = \omega^a(-1) \cdot p^f.
  \]
\end{lemma}

\begin{proof}
  \leanok
  Since $(q-1) \nmid a$, $\omega^{-a}$ is nontrivial and $\psi$ is primitive
  by \ref{lem:addCharTrace_isPrimitive}. By \texttt{gaussSum\_mul\_gaussSum\_eq\_card},
  $g(\omega^{-a}, \psi) \cdot g(\omega^{a}, \psi^{-1}) = p^f$.
  By the change of variable $x \mapsto -x$,
  $g(\omega^{a}, \psi^{-1}) = \omega^{a}(-1) \cdot g_{-a}$.
  Hence $g_a \cdot g_{-a} = \omega^a(-1) \cdot p^f$.
\end{proof}

\begin{lemma}
  \label{lem:norm_GaussSum}
  \uses{def:gauss_sum_s, lem:GaussSum_mul_GaussSum_neg, lem:GaussSum_mem}
  \lean{norm_GaussSum}
  \leanok
  If $(q-1) \nmid a$, then there exists $k \geq 1$ such that
  $N_{L/\QQ}(g_a) = \pm p^k$.
  In particular, the only rational prime dividing $N_{L/\QQ}(g_a)$ is $p$.
\end{lemma}

\begin{proof}
  \leanok
  By \ref{lem:GaussSum_mul_GaussSum_neg}, $g_a \cdot g_{-a} = \omega^a(-1) \cdot p^f = \pm p^f$.
  Taking norms, $N_{L/\QQ}(g_a) \cdot N_{L/\QQ}(g_{-a}) = \pm p^{f \cdot [L:\QQ]}$.
  Since both norms are integers and their product is $\pm p^{f \cdot [L:\QQ]}$,
  each is of the form $\pm p^k$ for some $k \geq 0$. By \ref{lem:GaussSum_mem},
  $g_a \in \fP$, so $g_a$ is not a unit in $\cO_L$, hence $N_{L/\QQ}(g_a) \neq \pm 1$,
  thus $k \geq 1$.
\end{proof}

\begin{lemma}
  \label{lem:mk_sq_gausssum_eq}
  \uses{lem:addCharTrace_mk_sq_eq, lem:orderOf_teichmuller}
  \lean{mk_sq_gausssum_eq}
  \leanok
  Assume $\zeta_p - 1 \in \fP$. Then
  \[
    g_1 \equiv -(\zeta_p - 1) \pmod{\fP^2}.
  \]
\end{lemma}

\begin{proof}
  \leanok
  By \ref{lem:addCharTrace_mk_sq_eq}, $\psi(x) \equiv 1 + \tr(x) \cdot (\zeta_p - 1) \pmod{\fP^2}$
  for all $x \in \cO_K/P$. Therefore
  \[
    g_1 \equiv \sum_{x} \omega(x)^{-1} + (\zeta_p-1) \sum_{x} \omega(x)^{-1} \tr(x) \pmod{\fP^2}.
  \]
  Since $p$ is odd, $q - 1 \geq 2$, so $\omega^{-1} \neq 1$ by
  \ref{lem:orderOf_teichmuller}, and $\sum_x \omega(x)^{-1} = 0$ by orthogonality
  of characters. Moreover,
  Computing the trace via powers of Frobenius, $\sum_{x \in \cO_K/P} \omega(x)^{-1} \cdot \tr(x) = -1$.
  Hence $g_1 \equiv -(\zeta_p - 1) \pmod{\fP^2}$.
\end{proof}

\begin{lemma}
  \label{lem:galFEquiv}
  \uses{def:gauss_sum_s}
  \lean{galFEquiv}
  \leanok
  The extension $L/F$ is Galois with $\Gal(L/F) \cong (\ZZ/(q-1)\ZZ)^\times$,
  via the action on $\zeta_{q-1}$.
\end{lemma}

\begin{proof}
  \leanok
  Since $\gcd(q-1, p) = 1$, the fields $K = \QQ(\zeta_{q-1})$ and
  $F = \QQ(\zeta_p)$ are linearly disjoint over $\QQ$. Therefore
  $L = KF$ is Galois over $F$ with $\Gal(L/F) \cong \Gal(K/\QQ) \cong
  (\ZZ/(q-1)\ZZ)^\times$.
\end{proof}

\begin{lemma}
  \label{lem:gal_gaussSum_eq_gaussSum}
  \uses{def:gauss_sum_s, lem:map_teichmuller_zpow_eq}
  \lean{gal_gaussSum_eq_gaussSum}
  \leanok
  For all $\sigma \in \Gal(L/F)$ and $a \in \ZZ$,
  \[
    \sigma(g_a) = g_{a \cdot \fn(\sigma)}.
  \]
\end{lemma}

\begin{proof}
  \leanok
  Since $\sigma \in \Gal(L/F)$, it fixes $\zeta_p$, hence $\psi$.
  By \ref{lem:map_teichmuller_zpow_eq}, $\sigma(\omega(x)) = \omega(x)^{\fn(\sigma)}$
  for all $x$. Therefore
  \[
    \sigma(g_a) = \sum_{x \in \cO_K/P} \sigma(\omega(x)^{-a}) \cdot \psi(x)
    = \sum_{x \in \cO_K/P} \omega(x)^{-a \cdot \fn(\sigma)} \cdot \psi(x)
    = g_{a \cdot \fn(\sigma)}.
  \]
\end{proof}

\begin{lemma}
  \label{lem:GaussSum_frob}
  \uses{def:gauss_sum_s, lem:addCharTrace_frob_apply}
  \lean{GaussSum_frob}
  \leanok
  For all $a \in \ZZ$, $g_{pa} = g_a$.
\end{lemma}

\begin{proof}
  \leanok
  The Frobenius map $\varphi : x \mapsto x^p$ is a bijection on $\cO_K/P$.
  Using the substitution $y = \varphi(x)$,
  \[
    g_{pa} = \sum_{x \in \cO_K/P} \omega(x)^{-pa} \psi(x)
    = \sum_{y \in \cO_K/P} \omega(\varphi^{-1}(y))^{-pa} \psi(\varphi^{-1}(y)).
  \]
  By multiplicativity of $\omega$, $\omega(\varphi^{-1}(y))^p = \omega(\varphi^{-1}(y)^p)
  = \omega(y)$, so $\omega(\varphi^{-1}(y))^{-pa} = \omega(y)^{-a}$.
  By \ref{lem:addCharTrace_frob_apply}, $\psi(\varphi^{-1}(y)) = \psi(y)$.
  Therefore $g_{pa} = g_a$.
\end{proof}

\section{Jacobi sums}

\begin{lemma}
  \label{lem:GaussSum_mul_GaussSum}
  \uses{def:gauss_sum_s}
  \lean{GaussSum_mul_GaussSum}
  \leanok
  If $(q-1) \nmid (a+b)$, then
  \[
    g_a \cdot g_b = g_{a+b} \cdot J(\omega^{-a}, \omega^{-b}).
  \]
\end{lemma}

\begin{proof}
  \leanok
  This follows from \texttt{jacobiSum\_mul\_nontrivial} applied to
  $\chi = \omega^{-a}$ and $\phi = \omega^{-b}$.
\end{proof}

\begin{lemma}
  \lean{JacobiSum}
  \leanok
  \label{lem:JacobiSum}
  \uses{def:gauss_sum_s}
  For all $a, b \in \ZZ$, $J(\omega^{-a}, \omega^{-b}) \in \cO_K$.
\end{lemma}

\begin{proof}
\leanok
\end{proof}

\chapter{The valuation function $s$}

\begin{definition}
  \lean{valGauss}
  \leanok
  \label{def:s}
  For $a \in \ZZ$, we define
  \[
    s(a) = v_{\fP}(g_a)
  \]
\end{definition}

\begin{lemma}
  \label{lem:valGauss_eq_zero}
  \uses{def:s, lem:gaussSum_one_left, lem:addCharTrace_ne_one, lem:orderOf_teichmuller}
  \lean{valGauss_eq_zero}
  \leanok
  If $(q - 1) \mid a$, then $s(a) = 0$.
\end{lemma}

\begin{proof}
  \leanok
  Since $(q-1) \mid a$, $\omega^{-a} = 1$ by \ref{lem:orderOf_teichmuller},
  so $g_a = \sum_{x} \psi(x)$, which equals $0$ by \ref{lem:gaussSum_one_left}
  since $\psi$ is nontrivial by \ref{lem:addCharTrace_ne_one}.
  Hence $s(a) = v_{\fP}(0) = 0$.
\end{proof}

\begin{lemma}
  \label{lem:valGauss_subadditive}
  \uses{def:s, lem:GaussSum_mul_GaussSum, lem:valGauss_eq_zero}
  \lean{valGauss_subadditive}
  \leanok
  For all $a, b \in \ZZ$, $s(a + b) \leq s(a) + s(b)$.
\end{lemma}

\begin{proof}
  \leanok
  If $(q-1) \mid (a+b)$ then $s(a+b) = 0$ by \ref{lem:valGauss_eq_zero} and the result is trivial.
  Otherwise, by \ref{lem:GaussSum_mul_GaussSum},
  \[
    g_a \cdot g_b = g_{a+b} \cdot J(\omega^{-a}, \omega^{-b}).
  \]
  Taking valuations and using the fact that $J(\omega^{-a}, \omega^{-b}) \in \cO_L$,
  hence $v_{\fP}(J) \geq 0$, we obtain $s(a) + s(b) \geq s(a+b)$.
\end{proof}

\begin{lemma}
  \label{lem:valGauss_add_valGauss_sub_self}
  \uses{def:s, lem:ramificationIdx_eq_p_sub_one, lem:valGauss_eq_zero, lem:GaussSum_periodic, lem:GaussSum_mul_GaussSum_neg}
  \lean{valGauss_add_valGauss_sub_self}
  \leanok
  Assume $(q-1) \mid k$. Then
  \[
    s(a) + s(k - a) = \begin{cases} 0 & \text{if } (q-1) \mid a \\
    (p-1) \cdot f & \text{otherwise.} \end{cases}
  \]
\end{lemma}

\begin{proof}
  \leanok
  Since $(q-1) \mid k$, we have $s(k-a) = s(-a)$ by periodicity of $g$:
  by \ref{lem:GaussSum_periodic}, $g_{a+(q-1)} = g_a$, so $s$ has period $q-1$.

  If $(q-1) \mid a$, then $(q-1) \mid (-a)$, so $s(a) = s(-a) = 0$ by \ref{lem:valGauss_eq_zero} and the result is clear.

  If $(q-1) \nmid a$, by \ref{lem:GaussSum_mul_GaussSum_neg},
  $g_a \cdot g_{-a} = \omega^{a}(-1) \cdot p^f$.
  Since $\omega^{a}(-1)$ is a root of unity, it is a unit in $\cO_L$, hence
  $v_{\fP}(\omega^{a}(-1)) = 0$. By \ref{lem:ramificationIdx_eq_p_sub_one},
  $v_{\fP}(p^f) = f \cdot e(\fP \mid p) = f(p-1)$.
  Therefore
  \[
    s(a) + s(-a) = v_{\fP}(g_a) + v_{\fP}(g_{-a}) = v_{\fP}(g_a \cdot g_{-a})
    = v_{\fP}(p^f) = (p-1) \cdot f.
  \]
\end{proof}

\begin{lemma}
  \label{lem:valGauss_frob}
  \uses{def:s, lem:GaussSum_frob}
  \lean{valGauss_frob}
  \leanok
  For all $a \in \ZZ$, $s(pa) = s(a)$.
\end{lemma}

\begin{proof}
  \leanok
  By \ref{lem:GaussSum_frob}, $g_{pa} = g_a$, hence
  $s(pa) = v_{\fP}(g_{pa}) = v_{\fP}(g_a) = s(a)$.
\end{proof}

\begin{lemma}
  \label{lem:valGauss_one}
  \uses{lem:GaussSum_mem, lem:mk_sq_gausssum_eq, lem:zeta_sub_one_not_mem_sq}
  \lean{valGauss_one}
  \leanok
  $s(1) = 1$.
\end{lemma}

\begin{proof}
  \leanok
  By \ref{lem:GaussSum_mem}, $g_1 \in \fP$, so $s(1) \geq 1$.
  By \ref{lem:mk_sq_gausssum_eq}, $g_1 \equiv -(\zeta_p - 1) \pmod{\fP^2}$.
  Since $\zeta_p - 1 \notin \fP^2$ by \ref{lem:zeta_sub_one_not_mem_sq},
  we have $g_1 \notin \fP^2$, so $s(1) \leq 1$.
\end{proof}

\begin{lemma}
  \label{lem:valGauss_le_self}
  \uses{def:s, lem:valGauss_subadditive, lem:valGauss_one, lem:valGauss_eq_zero}
  \lean{valGauss_le_self}
  \leanok
  For all $a \in \ZZ$ with $a \geq 0$, $s(a) \leq a$.
\end{lemma}

\begin{proof}
  \leanok
  By induction on $a$. The case $a = 0$ follows from \ref{lem:valGauss_eq_zero}
  (since $(q-1) \mid 0$). For $a \geq 1$, by \ref{lem:valGauss_subadditive} and
  the induction hypothesis,
  \[
    s(a) \leq s(1) + s(a-1) \leq 1 + (a-1) = a. \qedhere
  \]
\end{proof}

\begin{lemma}
  \label{lem:valGauss_toNat_le_sum_digits}
  \uses{def:s, lem:valGauss_subadditive, lem:valGauss_frob, lem:valGauss_le_self}
  \lean{valGauss_toNat_le_sum_digits}
  \leanok
  Let $0 \leq a$ and let
  $a = a_0 + a_1 p + \cdots + a_{t-1} p^{t-1}$ be the $p$-adic expansion of $a$.
  Then
  \[
    s(a) \leq a_0 + a_1 + \cdots + a_{t-1}.
  \]
\end{lemma}

\begin{proof}
  \leanok
  By induction on the list of digits $(a_0, \ldots, a_{t-1})$.
  The base case $a = 0$ is immediate. For the inductive step,
  write $a = a_0 + p \cdot a'$ where $a' = a_1 + a_2 p + \cdots + a_{t-1} p^{t-2}$.
  By \ref{lem:valGauss_subadditive},
  \[
    s(a_0 + p \cdot a') \leq s(a_0) + s(p \cdot a').
  \]
  By \ref{lem:valGauss_le_self}, $s(a_0) \leq a_0$.
  By \ref{lem:valGauss_frob}, $s(p \cdot a') = s(a')$.
  By the induction hypothesis, $s(a') \leq a_1 + \cdots + a_{t-1}$.
  Hence $s(a) \leq a_0 + a_1 + \cdots + a_{t-1}$.
\end{proof}

\begin{lemma}
  \label{lem:s_sum}
  \uses{def:s, lem:valGauss_add_valGauss_sub_self, lem:valGauss_eq_zero}
  \leanok
  We have
  \[
    \sum_{a=0}^{q-2} s(a) = \frac{(q-2) \cdot f \cdot (p-1)}{2}.
  \]
\end{lemma}

\begin{proof}
  \leanok
  We pair each $a \in \{0, \ldots, q-2\}$ with $q-1-a$. By \ref{lem:valGauss_add_valGauss_sub_self}
  with $k = q-1$, for any $a$ with $(q-1) \nmid a$,
  \[
    s(a) + s(q-1-a) = (p-1) \cdot f.
  \]
  Note that $s(0) = 0$ and $s(q-1) = 0$ by \ref{lem:valGauss_eq_zero}.

  Since $p$ is odd, $q - 1 = p^f - 1$ is even. The set $\{1, \ldots, q-2\}$
  consists of $(q-3)/2$ pairs $\{a, q-1-a\}$ with $a \neq q-1-a$, plus the
  central term $a = (q-1)/2$ which pairs with itself. Since $q \geq 3$, we have
  $1 \leq (q-1)/2 \leq q-2$, so $(q-1) \nmid (q-1)/2$, and
  \ref{lem:valGauss_add_valGauss_sub_self} gives $2s((q-1)/2) = (p-1)f$,
  hence $s((q-1)/2) = (p-1)f/2$. Therefore
  \[
    \sum_{a=0}^{q-2} s(a) = \frac{q-3}{2} \cdot (p-1)f + \frac{(p-1)f}{2}
    = \frac{(q-2)(p-1)f}{2}.
  \]
\end{proof}

\begin{proposition}
  \label{prop:valGauss_toNat_eq_sum_digits}
  \uses{lem:valGauss_toNat_le_sum_digits, lem:s_sum}
  \lean{valGauss_toNat_eq_sum_digits}
  \leanok
  Let $0 \leq a < q - 1$ and let
  $a = a_0 + a_1 p + \cdots + a_{f-1} p^{f-1}$, $0 \leq a_i \leq p - 1$,
  be the $p$-adic expansion of $a$. Then
  \[
    s(a) = a_0 + a_1 + \cdots + a_{f-1}.
  \]
\end{proposition}

\begin{proof}
  \leanok
  By \ref{lem:valGauss_toNat_le_sum_digits}, $s(a) \leq a_0 + \cdots + a_{f-1}$ for all $a$.
  Summing over $0 \leq a \leq q - 2$, we get
  \[
    \sum_{a=0}^{q-2} s(a) \leq \sum_{a=0}^{q-2} (a_0 + \cdots + a_{f-1}).
  \]
  By \texttt{Nat.sum\_sum\_digits\_eq}, the right-hand side equals
  $\frac{(q-2) \cdot f \cdot (p-1)}{2}$, which equals the left-hand side
  by \ref{lem:s_sum}. Since each term satisfies
  $s(a) \leq a_0 + \cdots + a_{f-1}$ and the sums are equal, we must have
  $s(a) = a_0 + \cdots + a_{f-1}$ for all $a$.
\end{proof}

\chapter{Ideal factorization of Gauss sums}

\section{Setup for Stickelberger's theorem}

Let $m$ be a positive integer dividing $q-1$ such that $f$ is the
multiplicative order of $p$ modulo $m$. Let $d = (q-1)/m$.

Let $E = \QQ(\zeta_m, \zeta_p)$. We denote by $\tilde{\fp}$ the prime
ideal of $\cO_E$ lying below $\fP$, i.e.\ $\tilde{\fp} = \fP \cap \cO_E$,
and by $\tilde{\fp}_0$ the prime ideal of $\cO_{\QQ(\zeta_m)}$ lying below
$\fP$, i.e.\ $\tilde{\fp}_0 = \fP \cap \cO_{\QQ(\zeta_m)}$.

\begin{lemma}
  \label{lem:gauss_sum_in_cyclotomic}
  \uses{def:gauss_sum_s}
  We have $g_d \in \cO_E$.
\end{lemma}

\begin{proof}
\end{proof}

Let $G = \Gal(\QQ(\zeta_m)/\QQ) \cong (\ZZ/m\ZZ)^\times$.
For $(a, m) = 1$, let $\sigma_a \in G$ denote the element with
$\sigma_a(\zeta_m) = \zeta_m^a$. The \emph{Stickelberger element} is
\[
  \Theta = \Theta(\QQ(\zeta_m)) = \sum_{\substack{a = 1 \\ (a,m) = 1}}^{m}
  a \cdot \sigma_a^{-1} \in \ZZ[G].
\]
\begin{remark}
  With the notation of Washington, $\Theta = m\theta$ where
  $\theta = \sum_{(a,m)=1} \{a/m\} \sigma_a^{-1} \in \QQ[G]$.
\end{remark}

\section{Descent and Stickelberger's theorem}

\begin{lemma}
  \label{lem:s_fractional_part}
  \uses{prop:valGauss_toNat_eq_sum_digits}
  For $0 \leq h < q - 1$,
  \[
    s(h) = (p-1) \sum_{i=0}^{f-1} \left\{ \frac{p^i h}{q-1} \right\}.
  \]
\end{lemma}

\begin{proof}
  Let $h = a_0 + a_1 p + \cdots + a_{f-1} p^{f-1}$ be the $p$-adic
  expansion of $h$. By \ref{prop:valGauss_toNat_eq_sum_digits}, $s(h) = a_0 + \cdots + a_{f-1}$.
  On the other hand, $p^i h \equiv a_0 p^i + \cdots + a_{f-1} p^{i+f-1}
  \pmod{q-1}$, so
  \[
    \left\{ \frac{p^i h}{q-1} \right\} = \frac{1}{q-1}(a_0 p^i + a_1 p^{i+1}
    + \cdots + a_{f-1} p^{i-1}),
  \]
  where the indices are taken mod $f$. Summing over $i = 0, \ldots, f-1$,
  each digit $a_j$ appears multiplied by $(p^0 + p^1 + \cdots + p^{f-1})/(q-1) = 1/(p-1)$.
  Therefore
  \[
    (p-1) \sum_{i=0}^{f-1} \left\{ \frac{p^i h}{q-1} \right\} = a_0 + \cdots + a_{f-1} = s(h).
  \]
\end{proof}

\begin{lemma}
  \label{lem:s_mod}
  \uses{lem:s_fractional_part}
  For $a \in (\ZZ/m\ZZ)^\times$,
  \[
    m \cdot s(ad) = (p-1) \sum_{i=0}^{f-1} (ap^i \bmod m).
  \]
\end{lemma}

\begin{proof}
  By \ref{lem:s_fractional_part} applied to $h = ad$,
  \[
    s(ad) = (p-1) \sum_{i=0}^{f-1} \left\{ \frac{p^i \cdot ad}{q-1} \right\}
           = (p-1) \sum_{i=0}^{f-1} \left\{ \frac{ap^i}{m} \right\}.
  \]
  Since $m \cdot \{ap^i/m\} = ap^i \bmod m$ for each $i$, multiplying by $m$
  gives the result.
\end{proof}

\begin{lemma}
  \label{lem:gauss_sum_power_fixed_1}
  \uses{lem:gal_gaussSum_eq_gaussSum}
  For any $\tau \in \Gal(L/\QQ(\zeta_m, \zeta_p))$,
  $\tau(g_d^m) = g_d^m$.
\end{lemma}

\begin{proof}
  Since $\tau$ fixes $\zeta_m$, we have $\fn(\tau) \equiv 1 \pmod{m}$, so
  $d \cdot \fn(\tau) \equiv d \pmod{q-1}$.
  By \ref{lem:gal_gaussSum_eq_gaussSum},
  $\tau(g_d) = g_{d \cdot \fn(\tau)} = g_d$.
  Therefore $\tau(g_d^m) = g_d^m$.
\end{proof}

\begin{lemma}
  \label{lem:gauss_sum_galois_action_K}
  \uses{def:gauss_sum_s, def:addchar}
  For any $\tau \in \Gal(L/K)$ with $\tau(\zeta_p) = \zeta_p^e$,
  \[
    \tau(g_d) = \omega^d(e)^{-1} \cdot g_d.
  \]
\end{lemma}

\begin{proof}
  Since $\tau \in \Gal(L/K)$, it fixes $K = \QQ(\zeta_{q-1})$, hence fixes
  all $(q-1)$-th roots of unity, hence fixes $\omega(x)$ for all $x$.
  Since $\tau(\zeta_p) = \zeta_p^e$, we have $\tau(\psi(x)) = \psi(ex)$. Therefore
  \[
    \tau(g_d) = \sum_{x \in \cO_K/P} \omega(x)^{-d} \psi(ex).
  \]
  By the substitution $y = ex$ (bijection on $\cO_K/P$ since
  $e \in (\ZZ/p\ZZ)^\times \subset \FF_q^\times$),
  \[
    \sum_x \omega(x)^{-d}\psi(ex) = \sum_y \omega(y/e)^{-d}\psi(y)
    = \omega^d(e)^{-1} \sum_y \omega(y)^{-d}\psi(y) = \omega^d(e)^{-1} g_d.
  \]
\end{proof}

\begin{lemma}
  \label{lem:gauss_sum_power_fixed_2}
  \uses{lem:gauss_sum_galois_action_K}
  For any $\tau \in \Gal(L/K)$,
  $\tau(g_d^m) = g_d^m$.
\end{lemma}

\begin{proof}
  By \ref{lem:gauss_sum_galois_action_K}, $\tau(g_d) = \omega^d(e)^{-1} g_d$
  where $e$ is such that $\tau(\zeta_p) = \zeta_p^e$.
  Since $md = q-1$, we have $\omega^d(e)^{-m} = \omega^{q-1}(e)^{-1} = 1$, hence
  $\tau(g_d^m) = \omega^d(e)^{-m} g_d^m = g_d^m$.
\end{proof}

\begin{lemma}
  \label{lem:gauss_sum_power_in_cyclotomic}
  \uses{def:gauss_sum_s, lem:gauss_sum_in_cyclotomic,
        lem:gauss_sum_power_fixed_1, lem:gauss_sum_power_fixed_2}
  We have $g_d^m \in \cO_{\QQ(\zeta_m)}$.
\end{lemma}

\begin{proof}
  By \ref{lem:gauss_sum_power_fixed_1}, $g_d^m$ is fixed by
  $\Gal(L/\QQ(\zeta_m, \zeta_p))$, so $g_d^m \in \QQ(\zeta_m, \zeta_p)$.
  By \ref{lem:gauss_sum_power_fixed_2}, $g_d^m$ is fixed by $\Gal(L/K)$.
  Since $\QQ(\zeta_m, \zeta_p) = KF$ and $\Gal(KF/\QQ(\zeta_m))$ is generated
  by the restrictions of $\Gal(L/\QQ(\zeta_m, \zeta_p))$ and $\Gal(L/K)$,
  we conclude $g_d^m \in \QQ(\zeta_m)$.
  Since $g_d \in \cO_L$, we have $g_d^m \in \cO_{\QQ(\zeta_m)}$.
\end{proof}

\begin{lemma}
  \label{lem:L_E_unramified}
  The extension $L/E$ is unramified above $p$.
\end{lemma}

\begin{proof}
  We have $L = E(\zeta_{q-1})$ and $E = \QQ(\zeta_m, \zeta_p)$. Since
  $\gcd(q-1, p) = 1$, the extension $\QQ(\zeta_{q-1})/\QQ$ is unramified
  at $p$. Since the ramification above $p$ in $L/\QQ$ comes entirely from
  the adjunction of $\zeta_p$, which is already in $E$, the extension $L/E$
  is unramified above $p$.
\end{proof}

\begin{lemma}
  \label{lem:p_totally_ramified_in_E}
  The extension $E/\QQ(\zeta_m)$ is totally ramified above $\tilde{\fp}_0$
  with ramification index $p - 1$. In particular,
  $\tilde{\fp}^{p-1} = \tilde{\fp}_0 \cdot \cO_E$.
\end{lemma}

\begin{proof}
  Since $\gcd(m, p) = 1$ (as $m \mid q - 1$ and $\gcd(q-1, p) = 1$),
  the fields $\QQ(\zeta_m)$ and $\QQ(\zeta_p)$ are linearly disjoint over
  $\QQ$, and $E = \QQ(\zeta_m) \cdot \QQ(\zeta_p)$. The extension
  $\QQ(\zeta_p)/\QQ$ is totally ramified at $p$ with ramification index
  $p - 1$. Since $\QQ(\zeta_m)/\QQ$ is unramified at $p$ (as $p \nmid m$),
  the extension $E/\QQ(\zeta_m)$ is totally ramified at $\tilde{\fp}_0$
  with ramification index $p - 1$.
\end{proof}

\begin{lemma}
  \label{lem:stabilizer_tP0}
  The stabilizer of $\tilde{\fp}_0$ in $G = \Gal(\QQ(\zeta_m)/\QQ)$ is
  the subgroup $\langle \sigma_p \rangle$, which has order $f$.
  In particular, for $a \in (\ZZ/m\ZZ)^\times$ and $0 \leq i \leq f-1$,
  \[
    \sigma_{ap^i}^{-1}(\tilde{\fp}_0) = \sigma_a^{-1}(\tilde{\fp}_0).
  \]
\end{lemma}

\begin{proof}
  Since $f$ is the multiplicative order of $p$ modulo $m$ by hypothesis,
  $\langle \sigma_p \rangle$ has order $f$ in $G$. The stabilizer of
  $\tilde{\fp}_0$ in $G$ is the decomposition group of $\tilde{\fp}_0$
  over $p$, which is generated by $\sigma_p$ since $p$ generates the
  decomposition group.
\end{proof}

\begin{lemma}
  \label{lem:valuation_gauss_sum_power}
  \uses{lem:gauss_sum_in_cyclotomic, lem:L_E_unramified,
        lem:p_totally_ramified_in_E,
        lem:gal_gaussSum_eq_gaussSum, lem:gauss_sum_galois_action_K, def:s,
        lem:s_mod}
  For $a \in (\ZZ/m\ZZ)^\times$,
  \[
    v_{\sigma_a^{-1}(\tilde{\fp}_0)}(g_d^m) = \sum_{i=0}^{f-1}
    (ap^i \bmod m).
  \]
\end{lemma}

\begin{proof}
  Since $g_d \in \cO_E$ by \ref{lem:gauss_sum_in_cyclotomic}, we have
  \[
    v_{\sigma_a^{-1}(\tilde{\fp})}(g_d) = v_{\tilde{\fp}}(\sigma_a(g_d)).
  \]
  Lift $\sigma_a \in G = \Gal(\QQ(\zeta_m)/\QQ)$ to an element
  $\tilde{\sigma}_a \in \Gal(L/F)$ with $\mathfrak{n}(\tilde{\sigma}_a) \equiv a
  \pmod{m}$. By \ref{lem:gal_gaussSum_eq_gaussSum},
  $\tilde{\sigma}_a(g_d) = g_{d \cdot \mathfrak{n}(\tilde{\sigma}_a)}$.
  The ambiguity in the lift is an element of $\Gal(L/\QQ(\zeta_m, \zeta_p))$,
  which acts on $g_d$ by \ref{lem:gauss_sum_galois_action_K} as multiplication
  by a unit $\omega^d(e)^{-1} \in \cO_E^\times$. Therefore
  $v_{\tilde{\fp}}(\sigma_a(g_d)) = v_{\tilde{\fp}}(g_{ad}) = s(ad)$,
  where the last equality uses that $L/E$ is unramified above $p$ by
  \ref{lem:L_E_unramified}, so $v_{\tilde{\fp}}(g_{ad}) = v_{\fP}(g_{ad}) = s(ad)$.
  Since $E/\QQ(\zeta_m)$ is totally ramified of degree $p-1$ above $\tilde{\fp}_0$
  by \ref{lem:p_totally_ramified_in_E}, we have
  $v_{\sigma_a^{-1}(\tilde{\fp})} = (p-1) \cdot v_{\sigma_a^{-1}(\tilde{\fp}_0)}$.
  Therefore
  \[
    v_{\sigma_a^{-1}(\tilde{\fp}_0)}(g_d^m) = \frac{m \cdot s(ad)}{p-1}
    = \sum_{i=0}^{f-1} (ap^i \bmod m),
  \]
  where the last equality is \ref{lem:s_mod}.
\end{proof}

\begin{lemma}
  \label{lem:valuation_stickelberger}
  \uses{lem:stabilizer_tP0}
  For $a \in (\ZZ/m\ZZ)^\times$,
  \[
    v_{\sigma_a^{-1}(\tilde{\fp}_0)}(\tilde{\fp}_0^{\Theta}) =
    \sum_{i=0}^{f-1} (ap^i \bmod m).
  \]
\end{lemma}

\begin{proof}
  By definition of $\Theta$ and \ref{lem:stabilizer_tP0}, regrouping the
  sum over $(\ZZ/m\ZZ)^\times$ by orbits under $\langle p \rangle$:
  \[
    v_{\sigma_a^{-1}(\tilde{\fp}_0)}(\tilde{\fp}_0^{\Theta})
    = \sum_{b \in (\ZZ/m\ZZ)^\times} b \cdot
      v_{\sigma_a^{-1}(\tilde{\fp}_0)}(\sigma_b^{-1}(\tilde{\fp}_0))
    = \sum_{i=0}^{f-1} (ap^i \bmod m),
  \]
  since $v_{\sigma_a^{-1}(\tilde{\fp}_0)}(\sigma_b^{-1}(\tilde{\fp}_0)) = 1$
  if $b \equiv ap^i \pmod{m}$ for some $i$, and $0$ otherwise.
\end{proof}

\begin{lemma}
  \label{lem:gauss_sum_factorization_cyclotomic}
  \lean{GaussSum_factorization}
  \uses{lem:gauss_sum_power_in_cyclotomic, lem:norm_GaussSum,
        lem:valuation_gauss_sum_power, lem:valuation_stickelberger,
        lem:stabilizer_tP0}
  \leanok
  \[
    (g_d^m) = \tilde{\fp}_0^{\Theta}
  \]
  as ideals of $\cO_{\QQ(\zeta_m)}$.
\end{lemma}

\begin{proof}
  \leanok
  By \ref{lem:gauss_sum_power_in_cyclotomic}, $g_d^m \in \cO_{\QQ(\zeta_m)}$.
  By \ref{lem:norm_GaussSum}, $g_d^m$ is only divisible by primes above $p$.
  Similarly, $\tilde{\fp}_0^{\Theta}$ only involves primes above $p$.
  It therefore suffices to check that for every prime $\mathfrak{q} =
  \sigma_a^{-1}(\tilde{\fp}_0)$ of $\cO_{\QQ(\zeta_m)}$ above $p$,
  \[
    v_{\mathfrak{q}}(g_d^m) = v_{\mathfrak{q}}(\tilde{\fp}_0^{\Theta}).
  \]
  By \ref{lem:valuation_gauss_sum_power},
  \[
    v_{\mathfrak{q}}(g_d^m) = \sum_{i=0}^{f-1} (ap^i \bmod m).
  \]
  By \ref{lem:valuation_stickelberger},
  \[
    v_{\mathfrak{q}}(\tilde{\fp}_0^{\Theta}) = \sum_{i=0}^{f-1} (ap^i \bmod m).
  \]
  The two sides are equal, which completes the proof.
\end{proof}

\begin{theorem}
  \label{thm:stickelberger}
  \uses{lem:gauss_sum_factorization_cyclotomic}
  The ideal $\tilde{\fp}_0^{\Theta}$ is principal in $\QQ(\zeta_m)$.
\end{theorem}

\begin{proof}
  By \ref{lem:gauss_sum_factorization_cyclotomic},
  $(g_d^m) = \tilde{\fp}_0^{\Theta}$.
\end{proof}

\chapter{The Kronecker-Weber theorem}

\section{Reduction to the prime degree case}

Throughout this chapter, $\zeta_m$ denotes a primitive $m$-th root of unity.
We say that an extension is \emph{unramified} if it is unramified at all
finite primes. A normal extension $K/F$ has \emph{exponent} $m$ if
$\Gal(K/F)$ has exponent $m$.

The goal of this chapter is to prove the following theorem.

\begin{theorem}
  \label{thm:kronecker_weber}
  \uses{lem:kw_reduce_to_prime_power, prop:kw_odd_prime_power, prop:kw_2_power}
  Every abelian extension of $\QQ$ is contained in a cyclotomic field.
\end{theorem}

\begin{proof}
\end{proof}

\begin{lemma}
  \label{lem:kw_abelian_cyclic_decomp}
  Every finite abelian extension $K/\QQ$ is the compositum of cyclic
  extensions of prime power degree.
\end{lemma}

\begin{proof}
  The Galois group $\Gal(K/\QQ)$ is a finite abelian group, hence a direct
  product of cyclic groups of prime power order. Each factor corresponds to
  a cyclic subextension of prime power degree, and $K$ is their compositum.
\end{proof}

\begin{lemma}
  \label{lem:kw_reduce_to_unramified_outside_p}
  Let $K/\QQ$ be a cyclic extension of prime power degree $p^m$. Then there
  exists a cyclic extension $F/\QQ$ of prime power degree $p^m$ unramified
  outside $p$ such that $K$ is cyclotomic if and only if $F$ is cyclotomic.
\end{lemma}

\begin{proof}
\end{proof}

% Strategy: if q \neq p is ramified in K/Q, the ramification group
% I_q \subseteq \Gal(K/Q) \cong \ZZ/p^m\ZZ is a nontrivial subgroup. One shows
% that there exists a cyclotomic extension L = Q(\zeta_{q^k}) such that KL = FL
% where F/Q is cyclic of degree p^m with q unramified. The key point is that the
% ramification of q in KL/L comes entirely from L, so can be absorbed.
% Iterating over all primes q \neq p ramified in K gives the result.
% See Marcus, Number Fields, Theorem 21.

\begin{lemma}
  \label{lem:kw_reduce_to_prime_power}
  \uses{lem:kw_abelian_cyclic_decomp, lem:kw_reduce_to_unramified_outside_p,
        prop:kw_odd_prime_power, prop:kw_2_power}
  It suffices to prove that every cyclic extension of $\QQ$ of prime power
  degree unramified outside $p$ is cyclotomic.
\end{lemma}

\begin{proof}
  By \ref{lem:kw_abelian_cyclic_decomp}, any finite abelian extension of
  $\QQ$ is a compositum of cyclic extensions of prime power degree. By
  \ref{lem:kw_reduce_to_unramified_outside_p}, each such extension can be
  replaced by a cyclotomic extension times a cyclic extension of prime power
  degree unramified outside $p$. By \ref{prop:kw_odd_prime_power} and
  \ref{prop:kw_2_power}, the latter is cyclotomic.
\end{proof}

\begin{lemma}
  \label{lem:kw_cyclic_compositum}
  If $K$ and $K'$ are two cyclic $p$-extensions of $\QQ$ and their
  compositum $KK'$ is cyclic, then $K \subseteq K'$ or $K' \subseteq K$.
\end{lemma}

\begin{proof}
  If $K \not\subseteq K'$ and $K' \not\subseteq K$, then $KK'$ contains
  a subgroup isomorphic to $(\ZZ/p\ZZ)^2$, contradicting the assumption
  that $\Gal(KK'/\QQ)$ is cyclic.
\end{proof}

\begin{lemma}
  \label{lem:kw_ramification_reduction}
  \uses{lem:kw_cyclic_compositum}
  Let $K/\QQ$ be a cyclic extension of prime power degree $p^m$. If
  $q \neq p$ is ramified in $K/\QQ$, then there exists a cyclic cyclotomic
  extension $L/\QQ$ such that $KL = FL$ for some cyclic extension $F/\QQ$
  of prime power degree in which $q$ is unramified.
\end{lemma}

\begin{proof}
\end{proof}

\begin{lemma}
  \label{lem:kw_minkowski}
  Every extension of $\QQ$ of degree greater than $1$ is ramified at some
  finite prime.
\end{lemma}

\begin{proof}
\end{proof}

% Proof uses Minkowski's theorem on discriminants: |d_K| > 1 for any
% nontrivial extension K/Q, so some prime divides d_K and is thus ramified.
% A Mathlib PR is in progress for this result.

\begin{proposition}
  \label{prop:kw_2_quadratic}
  The only quadratic extensions of $\QQ$ unramified outside $2$ are
  $\QQ(i)$, $\QQ(\sqrt{-2})$, and $\QQ(\sqrt{2})$. In particular, the
  maximal real abelian $2$-extension of $\QQ$ with exponent $2$ and
  unramified outside $2$ is $\QQ(\sqrt{2})$, which is the subfield of
  $\QQ(\zeta_8)$ of degree $2$.
\end{proposition}

\begin{proof}
  A quadratic extension of $\QQ$ unramified outside $2$ has discriminant
  a power of $2$, so it is of the form $\QQ(\sqrt{d})$ with $d \in \{-1, \pm 2\}$.
  This gives exactly the three fields listed. Among these, only $\QQ(\sqrt{2})$
  is real.
\end{proof}

\begin{proposition}
  \label{prop:kw_2_power}
  \uses{prop:kw_2_quadratic, lem:kw_cyclic_compositum, prop:kw_exponent_p}
  Let $K/\QQ$ be a cyclic extension of degree $2^m$ unramified outside $2$.
  Then $K$ is cyclotomic.
\end{proposition}

\begin{proof}
  If $m = 1$ we are done by \ref{prop:kw_2_quadratic}. If $m \geq 2$,
  assume first that $K$ is nonreal. Then $K(i)/K$ is a quadratic extension,
  and its maximal real subfield $M$ is cyclic of degree $2^m$ over $\QQ$.
  Since $K/\QQ$ is cyclotomic if and only if $M$ is, we may assume that
  $K$ is totally real.

  Let $K'$ be the maximal real subfield of $\QQ(\zeta_{2^{m+2}})$. If $K'K$
  is not cyclic, then it contains three real quadratic subfields unramified
  outside $2$, contradicting \ref{prop:kw_2_quadratic}. Thus $K'K$ is
  cyclic, and \ref{lem:kw_cyclic_compositum} implies $K = K'$.
\end{proof}

\begin{proposition}
  \label{prop:kw_odd_prime_power}
  \uses{prop:kw_exponent_p, lem:kw_cyclic_compositum, lem:kw_minkowski}
  Let $K/\QQ$ be a cyclic extension of odd prime power degree $p^m$
  unramified outside $p$. Then $K$ is cyclotomic.
\end{proposition}

\begin{proof}
  Let $K'$ be the subfield of degree $p^m$ in $\QQ(\zeta_{p^{m+1}})$.
  If $K'K$ is not cyclic, then it contains a subfield of type $(p, p)$
  unramified outside $p$, which contradicts \ref{prop:kw_exponent_p}.
  Thus $K'K$ is cyclic, and \ref{lem:kw_cyclic_compositum} implies
  $K = K'$.
\end{proof}

\section{The prime degree case}

\begin{proposition}
  \label{prop:kw_exponent_p}
  \uses{thm:stickelberger, lem:kw_kummer, lem:kw_abelian_kummer,
        lem:kw_split_prime, lem:kw_mu_pth_power_ideal, lem:kw_class_trivial,
        lem:kw_mu_unit, lem:kw_unit_root_of_unity, lem:kw_ramification_reduction}
  Let $p$ be an odd prime. The maximal abelian extension of $\QQ$ of
  exponent $p$ unramified outside $p$ is cyclic of degree $p$: it is
  the subfield of degree $p$ of $\QQ(\zeta_{p^2})$.
\end{proposition}

\begin{proof}
  Let $K/\QQ$ be a cyclic extension of prime degree $p$ unramified outside $p$.
  By the sequence of lemmas below, $L = KF = \QQ(\zeta_{p^2})$
  (see \ref{lem:kw_unit_root_of_unity}). Therefore $K \subseteq \QQ(\zeta_{p^2})$.
  Since $[K:\QQ] = p = [\QQ(\zeta_{p^2}):F]$ and $K \cap F = \QQ$, $K$ is
  the unique subfield of degree $p$ of $\QQ(\zeta_{p^2})$ linearly disjoint
  from $F$, which is the subfield of degree $p$ of $\QQ(\zeta_{p^2})$ over $\QQ$.
  Since any two cyclic extensions of degree $p$ unramified outside $p$ are
  contained in $\QQ(\zeta_{p^2})$, the maximal such extension is exactly this
  subfield.
\end{proof}

The proof of \ref{prop:kw_exponent_p} proceeds via the following lemmas.
Let $K/\QQ$ be a cyclic extension of prime degree $p$ unramified outside $p$.
Set $F = \QQ(\zeta_p)$ and $L = KF$.

\begin{lemma}
  \label{lem:kw_kummer}
  \uses{lem:kw_abelian_kummer}
  With the above notation, $L = F(\sqrt[p]{\mu})$ for some nonzero
  $\mu \in \cO_F$, and $L/F$ is unramified outside $p$.
\end{lemma}

\begin{proof}
  Since $[K:\QQ] = p$ and $[F:\QQ] = p-1$ are coprime, $K \cap F = \QQ$,
  so $[L:F] = [KF:F] = [K:\QQ] = p$. The extension $L/F$ is Galois with
  cyclic Galois group of order $p$, and $F$ contains the $p$-th roots of
  unity $\mu_p$. By Kummer theory (\texttt{IsCyclic} + \texttt{Kummer}),
  $L = F(\sqrt[p]{\mu})$ for some $\mu \in F^\times$. Multiplying by a
  $p$-th power if necessary, we may assume $\mu \in \cO_F$.

  For the ramification: let $q \neq p$ be a prime. Since $K/\QQ$ is
  unramified at $q$ and $F/\QQ$ is unramified at $q$ (as $F = \QQ(\zeta_p)$
  is only ramified at $p$), the compositum $L = KF$ is also unramified
  at $q$.
\end{proof}

\begin{lemma}
  \label{lem:kw_abelian_kummer}
  \uses{lem:kummer_abelian_criterion}
  The Kummer extension $L = F(\sqrt[p]{\mu})$ is abelian over $\QQ$ if
  and only if for every $\sigma_a \in \Gal(F/\QQ)$ there exists
  $\xi \in F^\times$ such that $\sigma_a(\mu) = \xi^p \mu^a$.
\end{lemma}

\begin{proof}
\end{proof}

% This is a general result on Kummer extensions; see Washington, Lemma 14.7
% or Hilbert, Satz 147. Moved to auxiliary results chapter; proof to be added later.

\begin{lemma}
  \label{lem:kw_split_prime}
  \uses{lem:kw_abelian_kummer}
  Let $\mathfrak{q}$ be a prime ideal of $F$ with $(\mu) = \mathfrak{q}^r \mathfrak{a}$
  with $\mathfrak{q} \nmid \mathfrak{a}$. If $p \nmid r$ and $L/\QQ$ is
  abelian, then $\mathfrak{q}$ splits completely in $F/\QQ$.
\end{lemma}

\begin{proof}
  Let $\sigma_a$ be an element of the decomposition group of $\mathfrak{q}$
  in $\Gal(F/\QQ)$, so $\sigma_a(\mathfrak{q}) = \mathfrak{q}$.
  Since $L/\QQ$ is abelian, by \ref{lem:kw_abelian_kummer} there exists
  $\xi \in F^\times$ such that $\sigma_a(\mu) = \xi^p \mu^a$. Since
  $\sigma_a(\mathfrak{q}) = \mathfrak{q}$, applying $\sigma_a$ to the
  factorization $(\mu) = \mathfrak{q}^r \mathfrak{a}$ gives
  $(\xi^p \mu^a) = \mathfrak{q}^r \sigma_a(\mathfrak{a})$, hence
  $\mathfrak{q}^r \| \xi^p \mu^a$. But also $(\mu^a) = \mathfrak{q}^{ar}\mathfrak{a}^a$,
  so $\mathfrak{q}^r \| \xi^p \mu^a$ implies $r \equiv ar \pmod{p}$.
  Since $p \nmid r$, this forces $a \equiv 1 \pmod{p}$,
  so $\sigma_a = 1$ and $\mathfrak{q}$ splits completely in $F/\QQ$.
\end{proof}

\begin{lemma}
  \label{lem:kw_mu_val_p}
  \uses{lem:kw_split_prime}
  For any prime $\mathfrak{q}$ of $\cO_F$, $p \mid v_{\mathfrak{q}}(\mu)$.
\end{lemma}

\begin{proof}
  Write $v_{\mathfrak{q}}(\mu) = r$. If $\mathfrak{q} \nmid p$, then
  $\mathfrak{q}$ does not split completely in $F/\QQ$ (since $F/\QQ$ is
  totally ramified at $p$ and unramified elsewhere, so primes above $q \neq p$
  have nontrivial decomposition group). By \ref{lem:kw_split_prime}, $p \mid r$.
  If $\mathfrak{q} = \mathfrak{p}_p = (1-\zeta_p)$, the same argument applies:
  $\mathfrak{p}_p$ does not split completely in $F/\QQ$, so $p \mid r$.
\end{proof}

\begin{lemma}
  \label{lem:kw_mu_pth_power_ideal}
  \uses{lem:kw_split_prime, lem:kw_abelian_kummer, lem:kw_mu_val_p}
  With the above notation, $(1 - \zeta_p) \nmid \mu$, and $(\mu) = \mathfrak{a}^p$
  for some ideal $\mathfrak{a}$ of $\cO_F$.
\end{lemma}

\begin{proof}
  By \ref{lem:kw_mu_val_p}, $p \mid v_{\mathfrak{q}}(\mu)$ for every prime
  $\mathfrak{q}$. Therefore $(\mu) = \mathfrak{a}^p$ where
  $\mathfrak{a} = \prod_{\mathfrak{q}} \mathfrak{q}^{v_{\mathfrak{q}}(\mu)/p}$.
  In particular, $v_{\mathfrak{p}_p}(\mu) \equiv 0 \pmod{p}$, so either
  $v_{\mathfrak{p}_p}(\mu) = 0$ (i.e. $(1-\zeta_p) \nmid \mu$) or
  $v_{\mathfrak{p}_p}(\mu) \geq p$. The case $v_{\mathfrak{p}_p}(\mu) \geq p$
  would mean $(1-\zeta_p)^p \mid \mu$, but since $L/F$ is unramified outside
  $p$ and $\mu$ generates the Kummer extension, replacing $\mu$ by
  $\mu/(1-\zeta_p)^p$ (which generates the same extension) we may assume
  $v_{\mathfrak{p}_p}(\mu) = 0$, i.e. $(1-\zeta_p) \nmid \mu$.
\end{proof}

\begin{lemma}
  \label{lem:kw_class_trivial}
  \uses{lem:kw_mu_pth_power_ideal, thm:stickelberger, lem:kw_abelian_kummer}
  The ideal class $[\mathfrak{a}]$ is trivial, i.e., $\mathfrak{a}$ is principal.
\end{lemma}

\begin{proof}
  From $(\mu) = \mathfrak{a}^p$ and $\sigma_a(\mu) = \xi^p \mu^a$
  (by \ref{lem:kw_abelian_kummer}), we get
  $\sigma_a(\mathfrak{a})^p = \sigma_a(\mathfrak{a}^p) = \sigma_a((\mu)) =
  (\xi^p \mu^a) = \mathfrak{a}^{pa}$. Since the ideal monoid is torsion-free,
  $\sigma_a(\mathfrak{a}) = \mathfrak{a}^a$, hence $\sigma_a([\mathfrak{a}]) =
  [\mathfrak{a}]^a$ for every $1 \leq a < p$.

  Now, $\mathfrak{a}$ is coprime to $p$ by \ref{lem:kw_mu_pth_power_ideal}.
  By \ref{thm:stickelberger} applied with $m = p$, the ideal
  $\mathfrak{a}^{\Theta}$ is principal in $\QQ(\zeta_p)$, where
  $\Theta = \sum_{a=1}^{p-1} a\sigma_a^{-1} \in \ZZ[G]$. We compute
  $[\mathfrak{a}]^{\Theta}$:
  \[
    [\mathfrak{a}]^{\Theta} = \prod_{a=1}^{p-1} \sigma_a^{-1}([\mathfrak{a}])^a
    = \prod_{a=1}^{p-1} [\mathfrak{a}]^{a^{-1} \cdot a} = [\mathfrak{a}]^{p-1},
  \]
  where we used $\sigma_a^{-1}([\mathfrak{a}]) = [\mathfrak{a}]^{a^{-1}}$
  (since $\sigma_{a^{-1}}([\mathfrak{a}]) = [\mathfrak{a}]^{a^{-1}}$).
  Since $(\mu) = \mathfrak{a}^p$, the ideal $\mathfrak{a}^p$ is principal,
  so $[\mathfrak{a}]^p = 1$ in the class group. Therefore
  $[\mathfrak{a}]^{p-1} = [\mathfrak{a}]^{-1}$, and since
  $[\mathfrak{a}]^{\Theta} = 1$ (as $\mathfrak{a}^{\Theta}$ is principal),
  we get $[\mathfrak{a}]^{-1} = 1$, i.e. $[\mathfrak{a}] = 1$.
\end{proof}

\begin{lemma}
  \label{lem:kw_mu_unit}
  \uses{lem:kw_class_trivial}
  There exists a unit $\eta \in \cO_F^\times$ such that $\mu = \alpha^p \eta$
  for some $\alpha \in \cO_F$.
\end{lemma}

\begin{proof}
  By \ref{lem:kw_class_trivial}, $\mathfrak{a} = (\alpha)$ for some
  $\alpha \in \cO_F$. Then $(\mu) = \mathfrak{a}^p = (\alpha^p)$, so
  $\mu/\alpha^p \in \cO_F^\times$. Setting $\eta = \mu/\alpha^p$ gives
  $\mu = \alpha^p \eta$.
\end{proof}

\begin{lemma}
  \label{lem:kw_unit_root_of_unity}
  \uses{lem:kw_mu_unit, lem:kw_abelian_kummer}
  The unit $\eta$ is a $p$-th power times a root of unity, so $\mu = \zeta_p^t$
  for some $t$. Therefore $L = \QQ(\zeta_{p^2})$.
\end{lemma}

\begin{proof}
  Write $\eta = \zeta_p^t \varepsilon$ where $\varepsilon$ is a unit in the
  maximal real subfield of $F$, so $\sigma_{-1}(\varepsilon) = \varepsilon$.
  Since $L/\QQ$ is abelian, by \ref{lem:kw_abelian_kummer} applied to
  $\sigma_{-1}$, there exists $\xi \in F^\times$ such that
  $\sigma_{-1}(\mu) = \xi^p \mu^{-1}$. Since $\mu = \alpha^p \zeta_p^t \varepsilon$,
  \[
    \sigma_{-1}(\mu) = \alpha^p \zeta_p^{-t} \varepsilon
    \quad \text{and} \quad
    \xi^p \mu^{-1} = \xi^p \alpha^{-p} \zeta_p^{-t} \varepsilon^{-1}.
  \]
  Equating, $\varepsilon^2 = (\xi/\alpha)^{-p}$, so $\varepsilon$ is a $p$-th
  power in $F^\times$. Therefore $\mu = \alpha^p \zeta_p^t \varepsilon =
  (\alpha')^p \zeta_p^t$ for some $\alpha' \in \cO_F$, and
  $L = F(\sqrt[p]{\mu}) = F(\sqrt[p]{\zeta_p^t}) \subseteq \QQ(\zeta_{p^2})$.
  Since $[L:F] = p = [\QQ(\zeta_{p^2}):F]$, we conclude $L = \QQ(\zeta_{p^2})$.
\end{proof}

\chapter{Auxiliary results}


\begin{lemma}
  \label{lem:kummer_abelian_criterion}
  Let $F$ be a number field containing $\mu_p$, let $\mu \in F^\times$, and
  let $L = F(\sqrt[p]{\mu})$. Let $G_0 \leq \Gal(F/\QQ)$ be a subgroup.
  Then $L/\QQ$ is Galois with $\Gal(L/\QQ)$ extending $G_0$ if and only if
  for every $\sigma \in G_0$ there exists $\xi \in F^\times$ such that
  $\sigma(\mu) = \xi^p \mu^{a(\sigma)}$ for some $a(\sigma) \in \ZZ$.
  In particular, $L/\QQ$ is abelian if and only if for every
  $\sigma_a \in \Gal(F/\QQ)$ there exists $\xi \in F^\times$ such that
  $\sigma_a(\mu) = \xi^p \mu^a$.
\end{lemma}

\begin{proof}
\end{proof}

% See Washington, Lemma 14.7 or Hilbert, Satz 147.

\begin{lemma}
  \label{lem:ramificationIdx_eq_p_sub_one}
  \lean{ramificationIdx_eq_p_sub_one}
  \leanok
  The ramification index of $\fP$ above $p$ is $p - 1$, i.e.,
  \[
    e(\fP \mid p) = p - 1.
  \]
\end{lemma}

\begin{proof}
  \leanok
  This follows from the general formula for the ramification index
  in cyclotomic extensions.
\end{proof}

\begin{lemma}
  \label{lem:zeta_sub_one_not_mem_sq}
  \uses{lem:ramificationIdx_eq_p_sub_one}
  \lean{zeta_sub_one_not_mem_sq}
  \leanok
  We have $\zeta_p - 1 \notin \fP^2$.
\end{lemma}

\begin{proof}
  \leanok
  By \ref{lem:ramificationIdx_eq_p_sub_one}, $e(\fP \mid p) = p - 1$. Since $\zeta_p$ is a
  primitive $p$-th root of unity, $(\zeta_p - 1)$ generates the unique prime
  above $p$ in $\QQ(\zeta_p)$, which has ramification index $p - 1$. Hence
  $v_{\fP}(\zeta_p - 1) = 1$, so $\zeta_p - 1 \notin \fP^2$.
\end{proof}

\begin{lemma}
  \label{lem:gaussSum_one_left}
  \uses{def:gauss_sum_s, def:addchar}
  \lean{gaussSum_one_left}
  \leanok
  For any additive character $\psi$ of $\cO_K/P$,
  \[
    g(1, \psi) = \begin{cases} q - 1 & \text{if } \psi = 1 \\ -1 & \text{otherwise,} \end{cases}
  \]
  where $\psi = 1$ denotes the trivial character (identically equal to $1$).
\end{lemma}

\begin{proof}
\leanok
\end{proof}
